\section*{Overview}

Dijkstra's algorithm is the shortest-path routine for graphs with nonnegative edge weights. It grows the known region
outward from a source, always finalizing the not-yet-finalized vertex with the smallest tentative distance.

This is one of the most reused graph patterns in contests. Once the priority-queue invariant is clear, the code is
short and reliable.

\section*{When to Use It}

Use it when:

\begin{itemize}
  \item all edge weights are nonnegative,
  \item you need distances from one source to many vertices,
  \item the graph is sparse enough for an adjacency list and heap to be natural,
  \item BFS is too weak because weights are not all equal.
\end{itemize}

If the edges are only \(0\) and \(1\), 0-1 BFS is often better. If negative edges exist, the key greedy invariant
breaks and Bellman-Ford or something stronger is needed.

\section*{Core Idea}

Maintain tentative distances \(\texttt{dist[v]}\). Initially only the source has distance \(0\); everyone else is
infinity.

The heap stores candidates \((d, v)\). Each step extracts the smallest candidate. If it matches the current stored
distance, then \(v\) is finalized and all outgoing edges are relaxed.

Relaxing \((v, to, w)\) means checking whether going through \(v\) improves \(\texttt{dist[to]}\):
\[
\texttt{dist[to]} > \texttt{dist[v]} + w.
\]

\section*{Key Insight}

The priority queue is not just a speed trick. It enforces the main correctness invariant:
\textbf{when a fresh state \((d, v)\) leaves the heap, \(d\) is already the true shortest distance to \(v\)}.

That is exactly why nonnegative edges matter. If an unseen path could later subtract weight, then the ``smallest
tentative distance is safe to finalize'' rule would no longer hold.

\section*{Operations / Main Technique}

The ordinary sparse-graph version uses:

\begin{itemize}
  \item an adjacency list,
  \item a min-heap,
  \item one distance array,
  \item optional parent pointers for restoring a shortest path.
\end{itemize}

\paragraph{Stale entries.}
I almost always use the lazy heap version: push improved states freely, and when popping, skip the state if it no
longer matches \(\texttt{dist[v]}\). That avoids a more complicated decrease-key structure.

\paragraph{Worked example.}
If the heap contains both \((9, x)\) and later \((5, x)\), then the first one is stale by the time it is popped.
Skipping it is not optional bookkeeping. It is how the lazy version stays correct and simple.

\section*{Correctness Intuition}

Assume the heap pops a fresh state \((d, v)\). Suppose for contradiction there were a strictly shorter path to \(v\).
Walk along that path from the source until you reach the first not-yet-finalized vertex. Its predecessor on the path
was already finalized earlier, so the relaxation from that predecessor would have inserted a candidate no larger than
the true path length.

That candidate should have been popped before \((d, v)\), contradicting the choice of \(v\) as the smallest fresh
tentative distance. Therefore \(d\) is final.

\section*{Complexity Analysis}

\begin{itemize}
  \item heap-based Dijkstra on a sparse graph: \(O((n + m)\log n)\),
  \item memory: \(O(n + m)\).
\end{itemize}

The heap receives at most one new state per successful relaxation, and each push/pop costs a logarithm.

\section*{Implementation}

The reference code keeps:

\begin{itemize}
  \item adjacency lists of \((to, weight)\),
  \item a \texttt{priority\_queue} with \texttt{greater<...>},
  \item \texttt{long long} distances,
  \item parent restoration for one concrete shortest path.
\end{itemize}

That is the version I trust most under contest pressure. Multi-source Dijkstra is just the same code with several
initial zero-distance pushes.

\section*{Common Pitfalls}

\begin{itemize}
  \item Running Dijkstra on graphs with negative edges.
  \item Forgetting to skip stale heap entries.
  \item Using \texttt{int} when path lengths need \texttt{long long}.
  \item Choosing infinity too small and overflowing during relaxation.
  \item Missing simpler special cases such as BFS or 0-1 BFS.
\end{itemize}

\section*{Variants / Extensions}

\begin{itemize}
  \item multi-source Dijkstra,
  \item Dijkstra on state graphs instead of ordinary vertices,
  \item counting shortest paths while relaxing,
  \item Dijkstra with potentials inside Johnson-style reweighting or min-cost-flow subroutines.
\end{itemize}

In CP, ``run Dijkstra on a product state'' is often more useful than the textbook graph-only framing.

\section*{Practice Problems}

\begin{itemize}
  \item CSES: \texttt{Shortest Routes I}.
  \item CSES: \texttt{Flight Routes}.
  \item Weighted shortest-path problems where the graph is built from problem states rather than explicit roads.
\end{itemize}

\section*{References}

\begin{itemize}
  \item \href{https://cp-algorithms.com/graph/dijkstra.html}{cp-algorithms: Dijkstra's algorithm}.
  \item \href{https://cp-algorithms.com/graph/dijkstra_sparse.html}{cp-algorithms: Dijkstra on sparse graphs}.
  \item \href{https://usaco.guide/CPH.pdf}{Competitive Programmer's Handbook}.
  \item \href{https://codeforces.com/catalog}{Codeforces Catalog}.
\end{itemize}
